%%%%%%%%%%%%%%%%%%%%%%%%%%%%%%%%%%%%%%%%%%%%%%%%%%%%%%%%%%%%%%%%%%%%%%
% Lớp 11 | Giới hạn. Hàm số liên tục | Giới hạn của dãy số
% Dạng 1D3?1-5: Cấp số nhân lùi vô hạn
% Số câu: 42  (sinh tự động bởi tach_cau.py; từ nay thầy sửa trực tiếp file này)
%%%%%%%%%%%%%%%%%%%%%%%%%%%%%%%%%%%%%%%%%%%%%%%%%%%%%%%%%%%%%%%%%%%%%%

%%%%%%%%%%%%%%%% TRẮC NGHIỆM 4 PHƯƠNG ÁN %%%%%%%%%%%%%%%%

% nguồn: Ok_Toan11_TracNghiem_Lan2
\begin{ex}%[1D3H1-5]%[ma:693daefe]
	 Cho số $a=3{,}13131313\ldots $ là số thập phân vô hạn tuần hoàn có chu kì là $13$, số $a$ được biểu diễn dưới dạng phân số tối giản dạng $a=\dfrac{x}{y}$, trong đó $x$ và $y$ là các số nguyên dương. Tìm tổng $x+y$.
	 \choice
		{\True  $x+y=490$}
		{$x+y=409$} 
		{$x+y=211$} 
		 {$x+y=130$}
	\loigiai{Ta có  $a=3+\dfrac{13}{100}+\dfrac{13}{100^2}+\dfrac{13}{100^3}+\cdots=3+\dfrac{\dfrac{13}{100}}{1-\dfrac{1}{100}}=\dfrac{310}{99}$.\\
		Suy  ra $x=310$, $y=99$.\\
		Vậy $x+y=409$. 
	 }
\end{ex}

% nguồn: Ok_Toan11_TracNghiem_Lan2
\begin{ex}%[1D3H1-5]%[ma:4de6085b]
	Tổng vô hạn $S=1+\dfrac{1}{3}+\dfrac{1}{3^2}+\cdots+\dfrac{1}{3^n}+\cdots$ có giá trị  bằng
	 \choice
		{$3$} 
		{\True $\dfrac{3}{2}$} 
		{$4$} 
		 {$2$}
	\loigiai{Ta có $1$; $\dfrac{1}{3}$; $\dfrac{1}{3^2}$;\ldots; $\dfrac{1}{3^n}$; \ldots là một cấp  số nhân lùi vô hạn với công bội $q=\dfrac{1}{3}<1$.\\
		Vậy 
		$S=1+\dfrac{1}{3}+\dfrac{1}{3^2}+\cdots+\dfrac{1}{3^n}+\cdots=\dfrac{1}{1-\dfrac{1}{3}}=\dfrac{3}{2}$. 
	 }
\end{ex}

% nguồn: Ok_Toan11_TracNghiem_Lan2
\begin{ex}%[1D3H1-5]%[ma:e4bf05bb]
	Tổng của cấp số nhân lùi vô hạn $\dfrac{1}{2}$; $-\dfrac{1}{4}$; $\ldots $; $\dfrac{(-1)^{n+1}}{2^{n}}$; $\ldots$  có giá trị bằng bao nhiêu?
	 \choice
		{\True  $\dfrac{1}{3}$}
		 {$1$}
		 {$-\dfrac{1}{3}$}
		{$-\dfrac{2}{3}$} 
	\loigiai{ 
		Cấp số nhân có công bội $q=-\dfrac{1}{2}$ và  $u_{1}=\dfrac{1}{2}$.\\
		Vậy $S=\dfrac{u_{1}}{1-q}=\dfrac{\dfrac{1}{2}}{1+\dfrac{1}{2}}=\dfrac{1}{3}$. 
	} 
\end{ex}

% nguồn: Ok_Toan11_TracNghiem_Lan2
\begin{ex}%[1D3H1-5]%[ma:642fa467]
	Tính  tổng $S=1+\dfrac{1}{2}+\dfrac{1}{4}+\dfrac{1}{8}+\cdots +\dfrac{1}{2^{n}}+\cdots$
	\choice 
		{\True $2$} 
		{$3$} 
		 {$1$}
		{$\dfrac{1}{2}$} 
	\loigiai{ 
		Ta có $S=1+\dfrac{1}{2}+\dfrac{1}{4}+\dfrac{1}{8}+\cdots  +\dfrac{1}{2^{n}}+\cdots =\dfrac{1}{1-\dfrac{1}{2}}=2$.
	 }
\end{ex}

% nguồn: Ok_Toan11_TracNghiem_Lan2
\begin{ex}%[1D3H1-5]%[ma:df83acd6]
	Tìm dạng hữu tỷ của số thập phân vô hạn tuần hoàn $P=2{,}13131313  \ldots$.
	\choice 
		{$P=\dfrac{212}{99}$} 
		 {$P=\dfrac{213}{100}$}
		 {$P=\dfrac{211}{100}$}
		{\True $P=\dfrac{211}{99}$} 
	\loigiai{ 
		Lấy máy tính bấm  từng phương án.
	} 
\end{ex}

% nguồn: Ok_Toan11_TracNghiem_Lan2
\begin{ex}%[1D3H1-5]%[ma:02d2b2c2]
	Tính tổng  $S=\dfrac{1}{2}-\dfrac{1}{6}+\dfrac{1}{18}-\cdots+\dfrac{(-1)^{n-1}}{2\cdot 3^{n-1}}+\cdots$.
	 \choice
		{$S=\dfrac{3}{4}$} 
		{$S=\dfrac{8}{3}$} 
		{$S=\dfrac{2}{3}$} 
		{\True $S=\dfrac{3}{8}$} 
	 \loigiai{
		Đây là tổng một cấp số nhân lùi vô hạn  có $\heva{&u_{1}=\dfrac{1}{2} \\& q=\dfrac{u_{2}}{u_{1}}=-\dfrac{1}{3}.}$\\
		 Áp dụng công thức $S=u_{1}+u_{2}+u_{3}+\ldots+u_{n}+\ldots=\dfrac{u_{1}}{1-q}\quad(|q|<1)$.\\
		 Tổng cần tính là $S=\dfrac{1}{2}-\dfrac{1}{6}+\dfrac{1}{18}-\cdots+\dfrac{(-1)^{n-1}}{2\cdot 3^{n-1}}+\cdots=\dfrac{\dfrac{1}{2}}{1+\dfrac{1}{3}}=\dfrac{3}{8}$.
	 }
\end{ex}

% nguồn: Ok_Toan11_TracNghiem_Lan2
\begin{ex}%[1D3H1-5]%[ma:dd3507f9]
	Tính tổng $S$ của cấp số nhân lùi vô  hạn có số hạng đầu $u_{1}=1$ và công bội $q=-\dfrac{1}{2}$.
	\choice 
		{$S=2$} 
		 {$S=\dfrac{3}{2}$}
		 {$S=1$}
		 {\True $S=\dfrac{2}{3}$}
	 \loigiai{
		 Ta có $S=\dfrac{u_{1}}{1-q}=\dfrac{1}{1-\left(-\dfrac{1}{2}\right)}=\dfrac{2}{3}$.
	} 
\end{ex}

% nguồn: Ok_Toan11_TracNghiem_Lan2
\begin{ex}%[1D3H1-5]%[ma:a9ebe02c]
	Tổng $S=1+\dfrac{1}{3}+\dfrac{1}{3^{2}}+\dfrac{1}{3^{3}}+\cdots +\dfrac{1}{3^{n}}+\cdots$ có giá  trị là
	 \choice
		{$-\dfrac{2}{3}$} 
		{\True  $\dfrac{3}{2}$}
		{$\dfrac{2}{3}$} 
		{$-\dfrac{3}{2}$} 
	\loigiai{ 
		Ta có $S$ là tổng  của cấp số nhân lùi vô hạn có $u_{1}=1$, $q=\dfrac{1}{3}$.\\
		 Suy ra $S=\dfrac{u_{1}}{1-q}=\dfrac{1}{1-\dfrac{1}{3}}=\dfrac{3}{2}$.
	} 
\end{ex}

% nguồn: Ok_Toan11_TracNghiem_Lan2
\begin{ex}%[1D3H1-5]%[ma:49582f6c]
	Cấp số nhân lùi vô hạn $\left(u_{n}\right)$ có $u_{1}=-2$; $q=\dfrac{1}{2}$. Khi đó tổng $S$ của cấp số  nhân đã cho bằng
	\choice 
		{$4$} 
		{$-\dfrac{4}{3}$} 
		 {\True $-4$}
		 {$\dfrac{4}{3}$}
	\loigiai{ 
		 Cấp số nhân lùi vô hạn $\left(u_{n}\right)$ có $u_{1}=-2$, $q=\dfrac{1}{2}$.\\
		Vậy $S=\dfrac{u_{1}}{1-q}=\dfrac{-2}{1-\dfrac{1}{2}}=-4$. 
	} 
\end{ex}

% nguồn: Ok_Toan11_TracNghiem_Lan2
\begin{ex}%[1D3H1-5]%[ma:a18bdb70]
	Cho tổng của một cấp số nhân lùi vô hạn $S=1-\dfrac{1}{3}+\dfrac{1}{9}-\dfrac{1}{27}+\dfrac{1}{81}-\cdots$. Giá  trị của $S$ là
	 \choice
		{$S=-\dfrac{3}{4}$} 
		{$S=-\dfrac{4}{3}$} 
		{\True  $S=\dfrac{3}{4}$}
		 {$S=\dfrac{4}{3}$}
	 \loigiai{
		Ta có dãy số  $1$, $-\dfrac{1}{3}$, $\dfrac{1}{9}$, $-\dfrac{1}{27}$, $\dfrac{1}{81}$, $\ldots$ là một cấp số nhân lùi vô hạn với $u_{1}=1$, $q=-\dfrac{1}{3}$.\\
		 Vậy $S=\dfrac{u_{1}}{1-q}=\dfrac{1}{1+\dfrac{1}{3}}=\dfrac{3}{4}$.
	} 
\end{ex}

% nguồn: Ok_Toan11_TracNghiem_Lan2
\begin{ex}%[1D3H1-5]%[ma:0b9b3429]
	 Tính tổng $S=16-8+4-2+\cdots$
	\choice 
		 {$32$}
		{\True $\dfrac{32}{3}$} 
		{$24$} 
		 {$-\dfrac{32}{3}$}
	\loigiai{ 
		Dãy số $16$, $-8$, $4$, $-2$, $\ldots$ là một cấp số nhân lùi vô hạn có số hạng đầu $u_{1}=16$ và  công bội $q=-\dfrac{1}{2}$.\\
		Vậy  $S=\dfrac{u_{1}}{1-q}=\dfrac{16}{1+\dfrac{1}{2}}=\dfrac{32}{3}$.
	 }
\end{ex}

% nguồn: Ok_Toan11_TracNghiem_Lan2
\begin{ex}%[1D3H1-5]%[ma:1aaf2fdd]
	Tam giác mà ba đỉnh của nó là ba trung điểm ba cạnh của tam giác $ABC$ được gọi là tam giác trung bình của tam giác $ABC$. Ta xây dựng dãy các tam giác $A_1 B_1 C_1$, $A_2 B_2 C_2$, $A_3 B_3 C_3$, $\ldots$ sao cho $A_1 B_1 C_1$ là một tam giác đều cạnh bằng $x$ và với mỗi số nguyên dương $n \geq 2$, tam giác $A_n B_n C_n$ là tam giác trung bình của tam giác $A_{n-1} B_{n-1} C_{n-1}$. Với mỗi số nguyên dương $n$, kí  hiệu $S_n$ tương ứng là diện tích tam giác $A_n B_n C_n$. Tính tổng $S=S_1+S_2+\cdots+S_n+\cdots$
	\choice 
		{\True $\dfrac{\sqrt{3}}{3}  x^2$}
		{$\sqrt{3}  x^2$}
		 {$\dfrac{\sqrt{3}}{2} x^2$}
		 {$\dfrac{2\sqrt{3}}{3} x^2$}
	 \loigiai{
		Với $n=1$ thì tam giác đều $A_1 B_1 C_1$ có cạnh bằng $x$ nên diện tích tam giác $A_1 B_1 C_1$ là $S_1=\dfrac{\sqrt{3}}{4}  x^2$.\\
		 Với $n=2$ thì tam giác đều $A_2 B_2 C_2$ có cạnh bằng $\dfrac{x}{2}$ nên diện tích tam giác $A_2 B_2 C_2$ là \[S_2=\dfrac{\sqrt{3}}{4}\left(\dfrac{x}{2}\right)^2.\]
		Như vậy tam giác đều $A_n B_n C_n$ có cạnh bằng $\dfrac{x}{2^{n-1}}$ nên diện tích tam giác $A_n B_n C_n$  là \[S_n=\dfrac{\sqrt{3}}{4}\left(\dfrac{x}{2^{n-1}}\right)^2.\]
		Khi đó ta được dãy $S_1$, $S_2$, $\ldots S_n \ldots$ là một cấp số nhân lùi vô hạn với  số hạng đầu $u_1=S_1=\dfrac{\sqrt{3}}{4} x^2$ và công bội $q=\dfrac{1}{4}$.\\
		Do đó tổng $S=S_1+S_2+\cdots+S_n+\cdots=\dfrac{u_1}{1-q}=\dfrac{\sqrt{3}}{3}  x^2$.
	} 
\end{ex}

% nguồn: Ok_Toan11_TracNghiem_Lan2
\begin{ex}%[1D3V1-5]%[ma:4b059a57]
	Tính tổng cấp số nhân lùi vô hạn $-\dfrac{1}{2}$,  $\dfrac{1}{4}$, $-\dfrac{1}{8}$, $\ldots$, $\dfrac{\left(-1\right)^2}{2^n}$, $\ldots$ là
	 \choice
		 {$-1$}
		{$\dfrac{1}{2}$} 
		 {$-\dfrac{1}{4}$}
		{\True $-\dfrac{1}{3}$} 
	\loigiai{ 
		Cấp số nhân  có $u_1=-\dfrac{1}{2}$ công bội $q=-\dfrac{1}{2}$ nên tổng của cấp số nhân lùi vô hạng là\\
		$\lim S_n=\lim\dfrac{u_1  \left(1-q^n \right)}{1-q}=\dfrac{u_1}{1-q}=-\dfrac{1}{3}$.
	} 
\end{ex}

% nguồn: Ok_TracNghiemLop11_Lan1
\begin{ex}%[1D3V1-5]%[ma:62883fd5]
	Biểu diễn số thập phân vô hạn tuần  hoàn $0,(3)$ dưới dạng phân số tối giản dạng $\dfrac{a}{b}$ với $a,b\in \mathbb{N}^*$. Tính $a+b$.
	 \choice
		 {$7$}
		 {$6$}
		{$5$} 
		 {\True $4$}
	 \loigiai{
		 Ta có $0,(3)=\dfrac{3}{10}+\dfrac{3}{{{10}^2}}+\cdots+\dfrac{3}{{{10}^n}}+\cdots=\dfrac{\dfrac{3}{10}}{1-\dfrac{1}{10}}=\dfrac{1}{3}$.\\
		Do đó  $a=1, b=3$ nên $a+b=4$.
	} 
\end{ex}

% nguồn: Ok_Toan11_TracNghiem_Lan2
\begin{ex}%[1D3C1-5]%[ma:602da415]
	Cho dãy số $(u_n)$ xác định bởi $u_1=3$, $2u_{n+1}=u_n+1$, $\forall n\geq1$. Gọi $S_n$ là  tổng $n$ số hạng đầu tiên của dãy số $(u_n)$. Tính $\lim S_n$.
	\choice 
		{$\lim S_n=-\infty$} 
		 {$\lim S_n=1$}
		{\True $\lim S_n=+\infty$} 
		 {$\lim S_n=-1$}
	\loigiai{Ta có $2u_{n+1}=u_n+1$, $\forall n\geq 1\Leftrightarrow u_{n+1}=\dfrac{1}{2}u_n+\dfrac{1}{2}$, $\forall n\geq 1$.\\ 
		 Đặt $v_n=u_n-1$, $\forall n\geq1$.
		Khi đó $v_{n+1}=u_{n+1}-1=\dfrac{1}{2}\left(u_n-1\right)=\dfrac{1}{2}v_n$,  $\forall n\geq 1$.\\
		Suy ra $(v_n)$  là một cấp số nhân có công bội $q=\dfrac{1}{2}$.\\
		Gọi $T_n$ là tổng $n$ số hạng  đầu tiên của $(v_n)$.\\
		 Ta có $T_n=v_1\dfrac{1-\left(\dfrac{1}{2}\right)^n}{1-\dfrac{1}{2}}=2v_1\left[1-\left(\dfrac{1}{2}\right)^n\right]\Rightarrow S_n=T_n+n\Rightarrow\lim S_n=+\infty$
	 }
\end{ex}

% nguồn: Ok_Toan11_TracNghiem_Lan2
\begin{ex}%[1D3C1-5]%[ma:f42673af]
	Tính tổng các nghiệm $x\in(0;\pi)$ của  phương trình $$\sin x+\sin^2x+\sin^3x+\cdots+\sin^nx+\cdots=1.$$
	\choice 
		 {$\dfrac{\pi}{6}$}
		{\True $\pi$} 
		{$\dfrac{5\pi}{6}$} 
		 {$\dfrac{2\pi}{3}$}
	\loigiai{Thay $x=\dfrac{\pi}{2}$ thay vào phương trình ta  có $1+1+1+..+1+\cdots=1$ vô lý.\\
		Với $x\in(0;\pi)\setminus\left\{\dfrac{\pi}{2}\right\}$ thì $\sin x$; $\sin^2x$; $\sin^3x$; $\sin^nx$;\ldots là cấp số nhân lùi vô  hạn công bội $q=\sin x$.\\
		Do đó, vế trái  của phương trình là tổng của cấp số nhân lùi vô hạn nên
		\allowdisplaybreaks\begin{eqnarray*}
		&&\sin x+\sin^2x+\sin^3x+\cdots+\sin^nx+\cdots=1\\ 
		&\Leftrightarrow  &\dfrac{\sin x}{1-\sin x}=1\\
		 &\Leftrightarrow &\sin x=\dfrac{1}{2}\\
		&\Leftrightarrow&\hoac{&x=\dfrac{\pi}{6}\\&x=\dfrac{5\pi}{6}}\quad  \left(\text{do } x\in (0;\pi)\setminus\left\{\dfrac{\pi}{2}\right\}\right).
	\end{eqnarray*}
	Vậy tổng các nghiệm $x\in(0;\pi)$  là $\dfrac{\pi}{6}+\dfrac{5\pi}{6}=\pi$.
 }
\end{ex}


%%%%%%%%%%%%%%%% ĐÚNG/SAI %%%%%%%%%%%%%%%%

% nguồn: Ok_dungsai_lop11_Lan1
\begin{ex}%[1D3N1-5]%[ma:57d4cdb3]
 Mỗi  mệnh đề dưới đây đúng hay sai?
  \choiceTF
  {\True  $\lim\limits_{n \to +\infty} q^n=0$, với $\left|q\right|<1$}
  {$\lim\limits_{n  \to +\infty} q^n=+\infty$, với $\left|q\right|>1$}
  {$1$; $\dfrac{1}{2}$; $\dfrac{1}{4}$; $\dfrac{1}{8}$ là cấp số nhân lùi vô  hạn}
   {\True Cấp số nhân lùi vô hạn có số hạng đầu $u_1=1$, công bội $q=-\dfrac{1}{3}$ có tổng bằng $\dfrac{3}{4}$}
  \loigiai{
  \begin{itemchoice}
   \itemch $\lim\limits_{n  \to +\infty} q^n=0$, với $\left|q\right|<1$.
   \itemch Với $q>1$  thì $\lim\limits_{n \to +\infty} q^n=+\infty$; với $q<-1$ thì không tồn tại $\lim\limits_{n \to +\infty} q^n$.
   \itemch Dãy số $1$;  $\dfrac{1}{2}$; $\dfrac{1}{4}$; $\dfrac{1}{8}$ có hữu hạn số hạng nên không phải là cấp số nhân lùi vô hạn.
   \itemch Tổng cấp số nhân  lùi vô hạn là $S=\dfrac{u_1}{1-q}=\dfrac{1}{1+\dfrac{1}{3}}=\dfrac{3}{4}$.
  \end{itemchoice}
 } 
\end{ex}

% nguồn: Ok_dungsai_lop11_Lan2
\begin{ex}%[1D3H1-5]%[ma:efc90b19]
	Tìm được tổng  của cấp số nhân lùi vô hạn sau $S=1+\dfrac{1}{2}+\dfrac{1}{4}+\dfrac{1}{8}+\cdots$ và $T=1-\dfrac{1}{3}+\dfrac{1}{3^2}-\cdots+(-1)^n \dfrac{1}{3^n}+\cdots$.
	 \choiceTF
		{\True $S=1+\dfrac{1}{2}+\dfrac{1}{4}+\dfrac{1}{8}+\ldots$ là tổng của cấp số  nhân lùi vô hạn có công bội $q=\dfrac{1}{2}$}
		{$T=1-\dfrac{1}{3}+\dfrac{1}{3^2}-\cdots+(-1)^n \dfrac{1}{3^n}+\ldots$ là tổng của cấp số nhân lùi vô hạn có công bội  $q=\dfrac{1}{3}$}
		{\True $S>  T$}
		 {$S=\dfrac{1}{T}$}
	\loigiai{ 
		\begin{itemchoice}
			\itemch Đây là tổng của cấp số nhân lùi vô hạn có số hạng đầu $u_1=1$, công  bội $q=\dfrac{1}{2}$.\\
			 Suy ra $S=1+\dfrac{1}{2}+\dfrac{1}{4}+\dfrac{1}{8}+\ldots+\left(\dfrac{1}{2}\right)^n+\ldots=\dfrac{1}{1-\dfrac{1}{2}}=2$.
			\itemch Đây  là tổng của cấp số nhân lùi vô hạn có số hạng đầu $u_1=1$, công bội $q=-\dfrac{1}{3}$.\\
			Vì  vậy $T=1-\dfrac{1}{3}+\dfrac{1}{3^2}-\ldots+\left(-\dfrac{1}{3}\right)^n+\cdots=\dfrac{1}{1+\dfrac{1}{3}}=\dfrac{3}{4}$.
			\itemch $S=2$, $T=\dfrac{3}{4} \Rightarrow S>  T$.
			 \itemch $S\neq\dfrac{1}{T}$.
		\end{itemchoice}
	 }
\end{ex}

% nguồn: Ok_dungsai_lop11_Lan2
\begin{ex}%[1D3H1-5]%[ma:7aa14338]
	Tìm  được tổng của cấp số nhân lùi vô hạn sau $S=1-\dfrac{1}{2}+\dfrac{1}{4}-\dfrac{1}{8}+\cdots$ và \break $T=1+\dfrac{1}{3}+\dfrac{1}{3^2}+\cdots+\dfrac{1}{3^n}+\cdots$ Khi đó
	\choiceTF 
		{\True $S=1-\dfrac{1}{2}+\dfrac{1}{4}-\dfrac{1}{8}+\cdots$ là tổng của cấp số nhân lùi vô hạn có công  bội $q=-\dfrac{1}{2}$}
		{\True $T=1+\dfrac{1}{3}+\dfrac{1}{3^2}+\cdots+\dfrac{1}{3^n}+\cdots$ là tổng của cấp số nhân  lùi vô hạn có công bội $q=\dfrac{1}{3}$}
		{$S>T$} 
		{\True $S=\dfrac{1}{T}$} 
	\loigiai{ 
		\begin{itemchoice}
			\itemch Đây là tổng của cấp số  nhân lùi vô hạn có số hạng đầu $u_1=1$, công bội
			 $q=-\dfrac{1}{2}$. \\
			Suy ra $S=1-\dfrac{1}{2}+\dfrac{1}{4}-\dfrac{1}{8}+\cdots +\left(-\dfrac{1}{2}\right)^n+\cdots =\dfrac{1}{1+\dfrac{1}{2}}=\dfrac{2}{3}$. 
			\itemch Đây là tổng của cấp số nhân lùi vô hạn  có số hạng đầu $u_1=1$, công bội
			$q=\dfrac{1}{3}$.\\ 
			Do đó, ta  có
			$T=1+\dfrac{1}{3}+\dfrac{1}{3^2}+\cdots+\dfrac{1}{3^n}+\cdots 
			 =\dfrac{1}{1-\dfrac{1}{3}}=\dfrac{3}{2}\cdot$
			 \itemch Từ kết quả trên, ta suy ra $S<T$.
			\itemch Từ kết quả trên,  ta suy ra $S=\dfrac{1}{T}$.
		\end{itemchoice}
	 }
\end{ex}

% nguồn: Ok_dungsai_lop11_Lan1
\begin{ex}%[1D3V1-5]%[ma:9d2f0664]
 Cho cấp số nhân lùi vô hạn $(u_n)$ có số hạng đầu bằng $2$ và công  bội bằng $\dfrac{1}{3}$.
 \choiceTF 
  {\True Số hạng thứ  hai của cấp số nhân là $u_2=\dfrac{2}{3}$}
  {Số  hạng thứ $5$ của cấp số nhân là $u_5=\dfrac{1}{81}$}
  {Tổng của $5$ số hạng đầu của cấp số nhân bằng  $\dfrac{242}{243}$}
  {\True Tổng  cấp số nhân lùi vô hạn $(u_n)$ bằng $3$}
 \loigiai{ 
 Ta  có $u_1 = 2$ và công bội $q = \dfrac{1}{3}$.
 \begin{itemchoice}
   \itemch $u_2=u_1\cdot q=2\cdot \dfrac{1}{3}=\dfrac{2}{3}$.
  \itemch  $u_5 = u_1q^4 = \dfrac{2}{3^4}= \dfrac{2}{81}$.
  \itemch Tổng của $5$ số hạng đầu của cấp số nhân là $S_5  = u_1\dfrac{1-q^5}{1-q} = 2\dfrac{1 - \left(\dfrac{1}{3}\right)^5}{1 - \dfrac{1}{3}} = 3\left[1 - \left(\dfrac{1}{3}\right)^5\right]=\dfrac{242}{81}$.
  \itemch Khi đó tổng của cấp số nhân lùi vô hạn đã  cho là
  $$S = \dfrac{u_1}{1-q} =\dfrac{2}{1-\dfrac{1}{3}}  =3.$$
 \end{itemchoice}
 } 
\end{ex}

% nguồn: Ok_dungsai_lop11_Lan1
\begin{ex}%[1D3V1-5]%[ma:c0ecc019]
 Viết được các số thập phân vô hạn tuần hoàn dưới dạng phân số tối giản, ta được: $0{,}212121\ldots=\dfrac{a}{b}$; $4{,}333\ldots=\dfrac{c}{d}$.  Khi đó:
  \choiceTF
  {\True $a+b=40$} 
  {Ba số $a$; $b$; $58$ tạo thành một cấp số cộng} 
   {$c+d=15$}
  {\True $\lim c=13$} 
 \loigiai{ 
   Ta có $0{,}212121\ldots=0{,}21+0{,}0021+0{,}000021+\cdots$\\
  Đây là tổng của cấp  số nhân lùi vô hạn với số hạng đầu $0,21$ và công bội $\dfrac{1}{100}$. Vì vậy $0{,}212121\ldots=0{,}21+0{,}0021+0{,}000021+\cdots=\dfrac{0{,}21}{1-\dfrac{1}{100}}=\dfrac{7}{33}$.\\
   Do đó $a=7, b=33$.\\
  Ta  có $0{,}333\ldots=0{,}3+0{,}03+0{,}003+\cdots$\\
  Đây là tổng của cấp  số nhân lùi vô hạn với số hạng đầu là $0,3$ và công bội là $\dfrac{1}{10}$.\\
  Vì vậy $4{,}333\ldots=4+0{,}3+0{,}03+0{,}003+\cdots=4+\dfrac{0{,}3}{1-\dfrac{1}{10}}=\dfrac{13}{3}$.\\ 
  Do đó $c=13,  d=3$.
  \begin{itemchoice}
   \itemch  Đúng. $a+b=7+33=40$.
   \itemch Sai. Ba số $a$; $b$; $58$ không  tạo thành một cấp số cộng vì $a+58 \ne 2b$.
   \itemch Sai. $c+d=16\ne15$. 
   \itemch Đúng. $\lim c=13$ vì $c=13$. 
  \end{itemchoice}
 } 
\end{ex}


%%%%%%%%%%%%%%%% TRẢ LỜI NGẮN %%%%%%%%%%%%%%%%

% nguồn: Ok_TraLoiNgan_Lop11_Lan1
\begin{ex}%[1D3H1-5]%[ma:1aabc498]
	Tìm tổng sau: ${S=\dfrac{1}{3}-\dfrac{1}{9}+\dfrac{1}{27}-\cdots+\dfrac{(-1)^{n+1}}{3^n}+\cdots}$ 
	\shortans{0,25} 
	\loigiai{ 
		Ta thấy ${S}$ là tổng của cấp số nhân lùi vô hạn với số hạng đầu ${u_1=\dfrac{1}{3}}$, công bội  ${q=-\dfrac{1}{3}}$.\\
		Vì vậy  ${S=\dfrac{u_1}{1-q}=\dfrac{\dfrac{1}{3}}{1+\dfrac{1}{3}}=\dfrac{1}{3} \cdot \dfrac{3}{4}=\dfrac{1}{4}}$.
	 }
\end{ex}

% nguồn: Ok_TraLoiNgan_Lop11_Lan2
\begin{ex}%[1D3H1-5]%[ma:8dafaf82]
	Tìm tổng của cấp số nhân lùi vô hạn sau  $S=1-\dfrac{1}{2}+\dfrac{1}{4}-\dfrac{1}{8}+\cdots$ Kết quả làm tròn đến hàng phần mười.
	\shortans{0,7} 
	 \loigiai{
		Ta cần áp  dụng công thức tổng cấp số nhân lùi vô hạn
		\[S=u_1+u_2+u_3+\cdots+u_n+\ldots=\dfrac{u_1}{1-q},\] trong đó $u_1$, $q$ theo thứ tự là số hạng đầu và công bội của  cấp số nhân đó.\\
		Đây là tổng của cấp số nhân lùi vô  hạn có số hạng đầu $u_1=1$, công bội $q=-\dfrac{1}{2}$ nên
		\begin{eqnarray*}
			 S&=&1-\dfrac{1}{2}+\dfrac{1}{4}-\dfrac{1}{8}+\ldots+\left(-\dfrac{1}{2}\right)^n+\ldots\\
			&=&\dfrac{1}{1+\dfrac{1}{2}}\\ 
			 &=&\dfrac{2}{3}.
		\end{eqnarray*}
	} 
\end{ex}

% nguồn: Ok_TraLoiNgan_Lop11_Lan2
\begin{ex}%[1D3H1-5]%[ma:47c853a1]
	Số thập phân vô hạn tuần hoàn $0{,}511111\ldots$ viết  dạng phân số có dạng $\dfrac{a}{b}$ với $a$; $b$ là các số tự nhiên và $\dfrac{a}{b}$ là phân số tối giản. Tính $|b-2a|$.
	\shortans{1} 
	\loigiai{ 
		Ta có $0{,}51111\ldots=0{,}5+0{,}01+0{,}001+0{,}0001+\cdots$\\ 
		Xét  tổng $0{,}01+0{,}001+0{,}0001+\cdots$\\
		Đây là tổng của  cấp số nhân lùi vô hạn với số hạng đầu là $u_1=0{,}01$ và công bội $q=\dfrac{1}{10}$.\\
		 Vì vậy
		\begin{eqnarray*}
			 0{,}51111\ldots&=&0{,}5+0{,}01+0{,}001+0{,}0001+\cdots\\
			 &=&0{,}5+\dfrac{u_1}{1-q}\\
			&=&0{,}5+\dfrac{0{,}01}{1-\dfrac{1}{10}}\\ 
			 &=&\dfrac{23}{45}.
		\end{eqnarray*}
		Suy  ra $a=23$, $b=45$.\\
		Vậy $|b-2a|=1$. 
	 }
\end{ex}

% nguồn: Ok_TraLoiNgan_Lop11_Lan2
\begin{ex}%[1D3H1-5]%[ma:4978d461]
	Tìm tổng  $S=\dfrac{1}{3}-\dfrac{1}{9}+\dfrac{1}{27}-\ldots+\dfrac{(-1)^{n+1}}{3^n}+\cdots$
	 \shortans{0,25}
	\loigiai{ 
		Ta thấy $S$ là tổng  của cấp số nhân lùi vô hạn với số hạng đầu $u_1=\dfrac{1}{3}$, công bội $q=-\dfrac{1}{3}$.\\
		Vì vậy 
		\begin{eqnarray*}
			 S&=&\dfrac{u_1}{1-q}\\
			&=&\dfrac{\dfrac{1}{3}}{1+\dfrac{1}{3}}\\ 
			 &=&\dfrac{1}{4}.
		\end{eqnarray*}
		Vậy $S=0{,}25$. 
	 }
\end{ex}

% nguồn: Ok_TraLoiNgan_Lop11_Lan2
\begin{ex}%[1D3H1-5]%[ma:e0f4c797]
	Biết số thập phân vô hạn tuần hoàn $0{,}212121\ldots=\dfrac{a}{b}$ (với $a$, $b$ là các số  dương có ước chung lớn nhất là $1$). Tính giá trị $a+b$.
	 \shortans{40}
	 \loigiai{
		Ta có 
		\begin{eqnarray*}
			 0{,}212121\ldots&=&0{,}21+0{,}0021+0{,}000021+\cdots\\
			&=&\dfrac{21}{100}+\dfrac{21}{100^2}+\dfrac{21}{100^3}+\cdots\\ 
			 &=&\dfrac{\dfrac{21}{100}}{1-\dfrac{1}{100}}\\
			&=&\dfrac{7}{33}. 
		\end{eqnarray*}
		 Do đó $a=7$, $b=33$.\\
		Vậy $a+b=40$. 
	} 
\end{ex}

% nguồn: Ok_TraLoiNgan_Lop11_Lan2
\begin{ex}%[1D3H1-5]%[ma:e54fda29]
	Số thập phân vô hạn tuần hoàn $0{,}271414\ldots$ viết dạng phân số  có dạng $\dfrac{m}{n}$ với $m$; $n$ là các số tự nhiên và $\dfrac{m}{n}$ là phân số tối giản. Tính $n-3m$.
	\shortans{1839} 
	 \loigiai{
		Ta có  $0{,}271414\ldots=0{,}27+0{,}0014+0{,}000014+\ldots$\\
		Xét tổng  $0{,}0014+0{,}000014+0{,}00000014+\ldots$\\
		Đây là tổng của một cấp số nhân lùi vô hạn có số hạng đầu là $u_1=0{,}0014$ và công bội  $q=\dfrac{1}{100}$.\\
		Vì vậy 
		\begin{eqnarray*}
			0{,}271414\ldots&=&0{,}27+\dfrac{u_1}{1-q}\\ 
			 &=&\dfrac{27}{100}+\dfrac{0{,}0014}{1-\dfrac{1}{100}}\\
			 &=&\dfrac{2\,687}{9\,900}.
		\end{eqnarray*}
		 Suy ra $m=2\,687$, $n=9\,900$.\\
		Vậy  $n-3m=1\,839$.
	 }
\end{ex}

% nguồn: Ok_TraLoiNgan_Lop11_Lan1
\begin{ex}%[1D3V1-5]%[ma:082d055b]
	Viết dạng phân số của số thập phân  vô hạn tuần hoàn ${0{,}511111\ldots}$ có dạng $\dfrac{a}{b}$ (với $a,b\in \mathbb{N}^*;\dfrac{a}{b}$ tối giản). Tính giá trị $a+b$.
	\shortans{68} 
	 \loigiai{
		Ta có ${0{,}51111\ldots=0{,}5+0{,}01+0{,}001+0{,}0001+\cdots}$\\ 
		Xét tổng ${0{,}01+0{,}001+0{,}0001+\cdots}$ Đây là tổng của cấp số nhân lùi vô hạn với số hạng đầu là  ${u_1=0{,}01}$ và công bội ${q=\dfrac{1}{10}}$. Vì vậy $$0{,}51111\ldots=0{,}5+0{,}01+0{,}001+0{,}0001+\cdots=0{,}5+\dfrac{u_1}{1-q}=0{,}5+\dfrac{0{,}01}{1-\dfrac{1}{10}}=\dfrac{23}{45}.
		$$ 
		 Do đó $a=23, b=45$ nên $a+b=68$.
	} 
\end{ex}

% nguồn: Ok_TraLoiNgan_Lop11_Lan1
\begin{ex}%[1D3V1-5]%[ma:1b19bdc4]
	Cho dãy số $(u_n)$, $n\in\mathbb{N^*}$ được viết dưới dạng khai triển là: $\dfrac{9}{10}$, $\dfrac{99}{100}$, $\ldots$,  $\dfrac{99...9}{10^n}$, $\ldots$\\
	Biết $\lim\limits_{n\to  +\infty}(u_1+u_2+\cdot+u_n-n)=a$, khi đó $9a+10$ bằng bao nhiêu?
	 \shortans{9}
	 \loigiai{
	 Ta có
	\begin{eqnarray*}
		 &&u_1+u_2+\cdot+u_n-n=(u_1-1)+(u_2-1)+\ldots+(u_n-1)\\
		 &=&-\left(\dfrac{1}{10}+\dfrac{1}{100}+\cdot+\dfrac{1}{10^2}\right) \\
		&=&-\dfrac{1}{10}\cdot \dfrac{1-\left(\dfrac{1}{10}\right)^n}{1-\dfrac{1}{10}}=-\dfrac{1}{9}\left[1-\left(\dfrac{1}{10}\right)^n\right]. 
	\end{eqnarray*}
	Do đó $\lim\limits_{n\to +\infty}(u_1+u_2+\cdot+u_n-n)=\lim\limits_{n\to  +\infty}-\dfrac{1}{9}\left[1-\left(\dfrac{1}{10}\right)^n\right]=-\dfrac{1}{9}$.\\
	 Từ đó ta có $a=-\dfrac{1}{9}$. Suy ra $9a+10=9$.
	 }
\end{ex}

% nguồn: Ok_TraLoiNgan_Lop11_Lan1
\begin{ex}%[1D3V1-5]%[ma:f5a6d1af]
	Viết dạng phân số của số thập phân vô hạn tuần hoàn ${0,271414\ldots}$ có dạng $\dfrac{a}{b}$ (với $a,b\in  \mathbb{N}^*;\dfrac{a}{b}$ tối giản). Tính giá trị $4a-b$.
	\shortans{848} 
	\loigiai{ 
		Ta  có ${0,271414\ldots=0,27+0,0014+0,000014+\cdots}$\\
		Xét tổng ${0,0014+0,000014+0,00000014+\cdots}$ Đây là tổng của một cấp  số nhân lùi vô hạn có số hạng đầu là ${u_1=0,0014}$ và công bội ${q=\dfrac{1}{100}}$.\\
		Vì vậy  $${0,271414\ldots=0,27+\dfrac{u_1}{1-q}=\dfrac{27}{100}+\dfrac{0,0014}{1-\dfrac{1}{100}}=\dfrac{2687}{9900}}.$$
		Do đó $a=2687, b=9900$ nên $4a-b=848$. 
	 }
\end{ex}


%%%%%%%%%%%%%%%% TỰ LUẬN %%%%%%%%%%%%%%%%

% nguồn: Ok_TuanLuan_Toan11_Lan2
\begin{ex}%[1D3H1-5]%[ma:1d37032d]
	Cho dãy  số $\left(u_n \right)$ được xác định bởi:
	$\heva{&u_1 = 1\\& u_{n+1} = u_{n} + \dfrac{1}{2^n}, \left(n \in  \mathbb{N}^*\right).}$
	 \loigiai{
		 Ta có
		 \[
		u_n = \left(u_n - u_{n-1} \right) + \left(u_{n-1} - u_{n-2} \right) + \cdots + \left(u_2 - u_1 \right) + u_1  = \left(\dfrac{1}{2} \right)^{n-1} + \left(\dfrac{1}{2} \right)^{n-2} + \cdots + \dfrac{1}{2} + 1.
		\] 
		Dãy  số $\left\{\left(\dfrac{1}{2} \right)^{n-1}, \left(\dfrac{1}{2} \right)^{n-2}, \ldots, \dfrac{1}{2} \right\}$ là một cấp số nhân có $n-1$ số hạng với số hạng đầu $u_1 = \dfrac{1}{2}$ và công bội $q = \dfrac{1}{2}$.
		Vậy 
		 \[
		u_n = S_{n-1} + 1 = \dfrac{1}{2} \cdot \dfrac{1  - {{\left(\dfrac{1}{2} \right)}^{n-1}}}{1 - \dfrac{1}{2}} + 1 = 2 - \left(\dfrac{1}{2} \right)^{n-1}.
		\] 
		Do  đó
		\[ 
		\lim \left(u_n - 2 \right) = \lim \left[-\left(\dfrac{1}{2} \right)^{n-1} \right] =  0.
		\] 
	 }
\end{ex}

% nguồn: Ok_TuanLuan_Toan11_Lan2
\begin{ex}%[1D3H1-5]%[ma:2cf7943e]
	Tìm số hạng $u_1$ của cấp số nhân lùi vô hạn $\left(u_n\right)$, biết $S=4$;  $q=\dfrac{1}{2}$.
	 \loigiai{
		Ta có $S=\dfrac{u_1}{1-q} \Rightarrow  4=\dfrac{u_1}{1-\dfrac{1}{2}} \Rightarrow u_1=2$.
	} 
\end{ex}

% nguồn: Ok_TuanLuan_Toan11_Lan2
\begin{ex}%[1D3H1-5]%[ma:753d3a89]
	Số thập phân vô hạn tuần hoàn $0{,}32111 \ldots$ được biểu diễn dưới dạng phân số tối  giản $\dfrac{a}{b}$, trong đó $a$, $b$ là các số nguyên dương. Tính $a-b$.
	 \loigiai{Ta có
		$0{,}32111  \ldots=\dfrac{32}{100}+\dfrac{1}{10^{3}}+\dfrac{1}{10^{4}}+\dfrac{1}{10^{5}}+\ldots$.\\
		Vì $\dfrac{1}{10^3}+\dfrac{1}{10^4}+\dfrac{1}{10^5}+\cdots$ là tổng  của cấp số nhân lùi vô hạn với số hạng đầu $u_1=\dfrac{1}{10^3}$, công bội $q=\dfrac{1}{10}$.\\
		 Nên $0{,}32111 \ldots =\dfrac{32}{100}+\dfrac{\dfrac{1}{10^{3}}}{1-\dfrac{1}{10}}=\dfrac{289}{900}$.\\
		Vậy  $a=289$, $b=900$. Do đó $a-b=289-900=-611$.
	} 
\end{ex}

% nguồn: Ok_TuanLuan_Toan11_Lan2
\begin{ex}%[1D3H1-5]%[ma:3a9b32c5]
	Tổng $S=1+0{,}9+(0{,}9)^2+(0{,}9)^3+\cdots+(0{,}9)^{n-1}+\cdots$ có  kết quả bằng bao nhiêu?
	\loigiai{ 
		 Ta có $S=1+0{,}9+(0{,}9)^2+(0{,}9)^3+\cdots+(0{,}9)^{n-1}+\cdots$ \\
		Nhận thấy $S$ là tổng của cấp số nhân lùi vô  hạng có $u_1=1$, $q=0{,}9$ nên $S=\dfrac{u_1}{1-q}=\dfrac{1}{1-0{,}9}=10$.\\
	 }
\end{ex}

% nguồn: Ok_TuanLuan_Toan11_Lan2
\begin{ex}%[1D3H1-5]%[ma:cd18f49d]
	Tìm  $\lim\left(\dfrac{1}{2}+\left(-\dfrac{1}{4}\right)+\dfrac{1}{8}+\cdots+\dfrac{\left(-1\right)^{n+1}}{2^n}\right)$.
	\loigiai{ 
		\begin{itemize}
			 \item \textbf{Cách 1.} $\lim\left(\dfrac{1}{2}+\left(-\dfrac{1}{4}\right)+\dfrac{1}{8}+\cdots+\dfrac{\left(-1\right)^{n+1}}{2^n}\right)=\lim\left[\dfrac{1}{2}\cdot\dfrac{1-\left(-\dfrac{1}{2}\right)^n}{1+\dfrac{1}{2}}\right]=\dfrac{1}{3}$.
			\item \textbf{Cách 2.} Sử dụng tổng của cấp số  nhân lùi vô hạn là $ \dfrac{1}{2}\cdot\dfrac{1}{1+\dfrac{1}{2}} =\dfrac{1}{3}$.
	\end{itemize}}
\end{ex}

% nguồn: Ok_TuanLuan_Toan11_Lan2
\begin{ex}%[1D3H1-5]%[ma:10a66ff9]
	Tính tổng của cấp số nhân lùi vô hạn $1,-\dfrac{1}{2}, \dfrac{1}{4},-\dfrac{1}{8}, \ldots,\left(-\dfrac{1}{2}\right)^{n-1}, \ldots$ 
	\loigiai{ 
		Theo đề cho ta có $u_1=1$,  $q=-\dfrac{1}{2}$ nên suy ra $S=\dfrac{u_1}{1-q}=\dfrac{1}{1+\dfrac{1}{2}}=\dfrac{2}{3}$.\\
	 }
\end{ex}

% nguồn: Ok_TuanLuan_Toan11_Lan2
\begin{ex}%[1D3H1-5]%[ma:28bbedcf]
	Tìm  $\lim\dfrac{1+\dfrac{1}{2}+\dfrac{1}{4}+\cdots+\dfrac{1}{2^n}}{1+\dfrac{1}{3}+\dfrac{1}{9}+\cdots+\dfrac{1}{3^n}}$.
	 \loigiai{
		\begin{itemize}
			\item  \textbf{Cách 1.} $\lim\dfrac{1+\dfrac{1}{2}+\dfrac{1}{4}+\cdots+\dfrac{1}{2^n}}{1+\dfrac{1}{3}+\dfrac{1}{9}+\cdots+\dfrac{1}{3^n}}=\lim\dfrac{1\cdot\dfrac{1-\left(\dfrac{1}{2}\right)^{n+1}}{1-\dfrac{1}{2}}}{1\cdot\dfrac{1-\left(\dfrac{1}{3}\right)^n}{1-\dfrac{1}{3}}}=\dfrac{4}{3}$.
			\item \textbf{Cách 2.} Sử  dụng tổng của cấp số nhân lùi vô hạn $ \dfrac{\dfrac{1}{1-\dfrac{1}{2}}}{\dfrac{1}{1-\dfrac{1}{3}}}=\dfrac{4}{3}$.
		\end{itemize}
	} 
\end{ex}

% nguồn: Ok_TuanLuan_Toan11_Lan2
\begin{ex}%[1D3H1-5]%[ma:88c878f4]
	Tìm $\lim\left(\dfrac{1}{5}+\dfrac{1}{5^2}+\dfrac{1}{5^3}+\cdots+\dfrac{1}{5^n}\right)$. 
	\loigiai{ 
		\begin{itemize}
			 \item \textbf{Cách 1.} $\lim\left(\dfrac{1}{5}+\dfrac{1}{5^2}+\dfrac{1}{5^3}+\cdots+\dfrac{1}{5^n}\right)=\lim\dfrac{1}{5}\cdot \dfrac{1-\left(\dfrac{1}{5}\right)^n}{1-\dfrac{1}{5}}=\dfrac{1}{4}$.
			\item \textbf{Cách  2.} Sử dụng tổng của cấp số nhân lùi vô hạn là $ \dfrac{1}{5}\cdot \dfrac{1}{1-\dfrac{1}{5}}=\dfrac{1}{4} $.
	\end{itemize}}
\end{ex}

% nguồn: Ok_TuanLuan_Toan11_Lan2
\begin{ex}%[1D3H1-5]%[ma:0447f631]
	Cho số thập phân vô hạn tuần hoàn $a=2{,}151515 \ldots$ (chu kỳ 15), $a$ được biểu diễn dưới dạng phân số tối giản, trong đó $m$, $n$ là các số nguyên dương. Tìm tổng $m+n$. 
	\loigiai{Ta có $a=2{,}151515  \ldots=2+\dfrac{15}{100}+\dfrac{15}{100^{2}}+\dfrac{15}{100^{3}}+\cdots$\\
		 Vì $\dfrac{15}{100}+\dfrac{15}{100^{2}}+\dfrac{15}{100^{3}}+\cdots$ là tổng của cấp số nhân lùi vô hạn với số hạng đầu $u_{1}=\dfrac{15}{100}$, công bội $q=\dfrac{1}{100}$.\\
		Nên $a=2+\dfrac{\dfrac{15}{100}}{1-\dfrac{1}{100}}=\dfrac{71}{33}$.\\ 
		Vậy $m=71$, $n=33$ nên $m+n=104$. 
	 }
\end{ex}

% nguồn: Ok_TuanLuan_Toan11_Lan2
\begin{ex}%[1D3H1-5]%[ma:cfae156e]
	 Cho số thập phân vô hạn tuần hoàn $a=0,212121\ldots$ (chu kỳ là $21$). Viết $a$ dưới dạng phân số.
	\loigiai{ 
		 \textit{Cách 1}: Giải bằng tự luận\\
		Ta có $a=0{,}212121\ldots=0{,}21+0{,}0021+0{,}000021+\cdots=21\left(\dfrac{1}{10^2}+\dfrac{1}{10^4}+\dfrac{1}{10^6}+\cdots\right)$ \\ 
		Tổng $S=\dfrac{1}{10^2}+\dfrac{1}{10^4}+\dfrac{1}{10^6}+\cdots$ là tổng cấp số nhân lùi vô hạn có  $u_1=\dfrac{1}{10^2}, q=\dfrac{1}{10^2}$.\\
		Khi đó:  $S=\dfrac{u_1}{1-q}=\dfrac{\dfrac{1}{10^2}}{1-\dfrac{1}{10^2}}=\dfrac{1}{99}$. Do đó $A=21\cdot \dfrac{1}{99}=\dfrac{7}{33}$.\\
		 \textit{Cách 2}: Giải nhanh bằng máy tính\\
		Nhập vào màn hình $0{,}(21)$  và ấn phím $=$ ta được kết quả $\dfrac{7}{33}$.
	} 
\end{ex}

% nguồn: Ok_TuanLuan_Toan11_Lan2
\begin{ex}%[1D3V1-5]%[ma:87354732]
	Tìm $\lim\dfrac{1+2+2^2+2^3+\cdots+2^n}{1+3+3^2+3^3+\cdots+3^n}$. 
	\loigiai{ 
	 $\lim\dfrac{1+2+2^2+\cdots+2^n}{1+3+3^2+\cdots+3^n}=\lim\dfrac{1\cdot\dfrac{1-2^{n+1}}{1-2}}{1\cdot\dfrac{1-3^{n+1}}{1-3}}=\lim\dfrac{2\cdot\left(1-2^{n+1}\right)}{1-3^{n+1}}\lim\dfrac{2\cdot\left[\left(\dfrac{1}{3}\right)^{n+1}-\left(\dfrac{2}{3}\right)^{n+1}\right]}{\left(\dfrac{1}{3}\right)^{n+1}-1}=0$.}
\end{ex}

% nguồn: Ok_TuanLuan_Toan11_Lan2
\begin{ex}%[1D3V1-5]%[ma:5dc18257]
	Cho $S=1+q+q^2+q^3+\cdots$, $|q|<1$; $T=1+Q+Q^2+Q^3+\cdots$, $|Q|<1$; $E=1+qQ+q^2Q^2+q^3Q^3+\cdots$.  Biểu thị biểu thức $E$ theo $S$, $T$.
	 \loigiai{
		Ta có $S=1+q+q^2+q^3+\cdots,|q|< 1$ là tổng  của cấp số nhân lùi vô hạn, có $u_1=1, q=q$.\\
		Khi đó $S=\dfrac{u_1}{1-q}=\dfrac{1}{1-q} \Rightarrow  q=\dfrac{S-1}{S}$.\hfill(1) \\
		Tương tự  $T=\dfrac{1}{1-Q} \Rightarrow Q=\dfrac{T-1}{T}$.\hfill(2)\\
		Lại có $E=1+q \cdot Q+q^2\cdot Q^2+q^3\cdot Q^3+\cdots$ là tổng của cấp số nhân lùi vô hạng công bội $q Q$ (vì $|q Q|< 1$,  và $u_1=1$) nên
		 \[E=\dfrac{1}{1-q Q}. \tag{3} \]
		Thay  (1), (2) vào (3), ta được
		\[E=\dfrac{1}{1-\dfrac{T-1}{T}  \cdot \dfrac{S-1}{S}}=\dfrac{ST}{S+T-1}. \]
	 }
\end{ex}
